\ChapterImagePrelim[cap:glosario]{Glosario}{./images/fondo.png}\label{cap:glosario}
\mbox{}\\
En este apartado se encuentran términos clave y conceptos relevantes utilizados a lo largo de este proyecto.

\section*{A}
\begin{description}
	\item[Activo de información:] Cualquier información o elemento relacionado con el tratamiento de la información que tenga valor para la organización, tales como sistemas, servicios, soportes, infraestructura tecnológica o recursos humanos asociados~\cite{ISO27000Glossary}.
	
	\item[Amenaza:] Causa potencial de un incidente no deseado que puede ocasionar daños a un sistema, servicio, infraestructura o a la organización~\cite{ISO27000Glossary}.
	
	\item[Arquitectura empresarial:] Enfoque orientado a describir, estructurar y relacionar los procesos, servicios, aplicaciones, tecnologías y actores de una organización para apoyar sus objetivos estratégicos y operacionales~\cite{togaf-standard}.
	
	\item[Arquitectura de seguridad:] Estructura física y lógica compuesta por procesos, controles, tecnologías y relaciones que permiten proteger los activos de información de una organización mediante la definición de dominios, mecanismos y políticas de seguridad~\cite{NISTSecurityArchitecture}.
\end{description}

\section*{C}
\begin{description}
	\item[Computación en la nube:] \textbf{Computación en la nube}: Modelo que permite el acceso bajo demanda, a través de la red, a un conjunto compartido de recursos computacionales configurables, tales como redes, servidores, almacenamiento, aplicaciones y servicios, los cuales pueden ser aprovisionados y liberados rápidamente con un esfuerzo mínimo de gestión o interacción con el proveedor del servicio~\cite{mell2011nist}.
	
	\item[Ciberseguridad:] Conjunto de procesos, tecnologías, controles y prácticas orientadas a proteger sistemas, redes, dispositivos, servicios e información frente a amenazas, accesos no autorizados, ataques o incidentes que puedan comprometer la confidencialidad, integridad y disponibilidad de la información~\cite{NISTCybersecurity}.

	\item[Control de seguridad:] Medida, política, proceso, práctica o mecanismo implementado con el propósito de modificar, reducir o gestionar riesgos relacionados con la seguridad de la información~\cite{ISO27000Glossary}.
	
	\item[Correlación de eventos:]Proceso de relacionar y analizar eventos provenientes de diferentes fuentes con el propósito de identificar comportamientos, anomalías o incidentes de seguridad dentro de una infraestructura~\cite{IBM_SIEM}.
\end{description}

\section*{D}
\begin{description}
	\item[Defensa en profundidad:] Estrategia de seguridad basada en la implementación de múltiples capas de protección, controles y mecanismos de seguridad distribuidos a través de la infraestructura con el fin de reducir riesgos y limitar el impacto de posibles amenazas o incidentes.
\end{description}

\section*{E}
\begin{description}
	\item[Evento de seguridad de la información:] Ocurrencia identificada en un sistema, servicio o red que puede indicar una violación de la política de seguridad, una falla de los controles implementados o una situación relevante para la seguridad de la información~\cite{ISO27000Glossary}.
	
	\item[Evaluación de riesgos:] Proceso de identificación, análisis y estimación de los riesgos que pueden afectar los activos, servicios o procesos de una organización~\cite{ISO27000Glossary}.
	
	\item[Estudio de mapeo sistemático (\SMS):] Un estudio de mapeo sistemático es el proceso de identificar, categorizar y analizar la bibliografía existente relevante para un determinado tema de investigación. El objetivo de un \SMS~es obtener una visión global de un tema de investigación concreto, presentar una evaluación imparcial de la bibliografía actual, identificar las lagunas en la investigación y recopilar pruebas para futuras investigaciones \citep{Salama2017}.
\end{description}

\section*{F}
\begin{description}
	\item[Firewall:]  Dispositivo o mecanismo de seguridad encargado de supervisar, filtrar y controlar el tráfico de red entrante y saliente de acuerdo con reglas y políticas previamente definidas, con el propósito de proteger sistemas, servicios e infraestructuras frente a accesos o comunicaciones no autorizadas~\cite{NISTFirewall}.
\end{description}


\section*{G}
\begin{description}
	\item[Gobernanza de seguridad:] Conjunto de procesos, políticas y mecanismos orientados a dirigir, supervisar y controlar la gestión de la seguridad de la información dentro de una organización~\cite{MicrosoftSecurityGovernance}.
	
	\item[Gestión de incidentes:] Conjunto de procesos orientados a detectar, reportar, evaluar, responder y tratar incidentes de seguridad de la información, así como a generar aprendizaje a partir de estos~\cite{ISO27000Glossary}.
\end{description}

\section*{I}
\begin{description}
	\item[Incidente de seguridad de la información] Evento o conjunto de eventos de seguridad inesperados o no deseados que pueden comprometer las operaciones de la organización y afectar la confidencialidad, integridad o disponibilidad de la información~\cite{ISO27000Glossary}.
	
	\item[ISO/IEC 27001:] Norma internacional que establece los requisitos para implementar, mantener y mejorar continuamente un Sistema de Gestión de Seguridad de la Información (SGSI) orientado a la protección de los activos de información de una organización~\cite{ISO27000Glossary}.
	
	\item[ISO/IEC 27017:] Norma internacional que proporciona directrices y controles de seguridad de la información específicos para servicios de computación en la nube, complementando los controles definidos en la ISO/IEC 27002~\cite{iso27017}.
\end{description}

\section*{M}
\begin{description}
	\item[Monitoreo:] Proceso de supervisar, verificar u observar continuamente el estado de un sistema, servicio, proceso o actividad con el fin de identificar comportamientos, cambios o eventos relevantes~\cite{ISO27000Glossary}.
\end{description}

\section*{O}
\begin{description}
	\item[Organización:]  Persona o conjunto de personas con funciones, responsabilidades, autoridades y relaciones definidas, orientadas al cumplimiento de objetivos específicos~\cite{ISO27000Glossary}.
\end{description}

\section*{P}
\begin{description}
	\item[Política de seguridad de la información:] Conjunto de directrices, reglas y lineamientos establecidos para regular la protección, gestión y uso seguro de los activos de información y recursos tecnológicos dentro de una infraestructura u organización~\cite{iso27001}.
\end{description}


\section*{R}
\begin{description}
	\item[Respuesta activa:] Mecanismo automatizado o semiautomatizado orientado a ejecutar acciones de contención, bloqueo o mitigación frente a eventos o incidentes detectados dentro de una infraestructura~\cite{WazuhActiveResponse}.
	
	\item[Riesgo:] Efecto de la incertidumbre sobre los objetivos de una organización. En el contexto de la seguridad de la información, el riesgo se asocia con la posibilidad de que amenazas exploten vulnerabilidades presentes en los activos de información, generando impactos sobre la confidencialidad, integridad o disponibilidad de la información~\cite{ISO27000Glossary}.
\end{description}

\section*{S}
\begin{description}
	\item[\textit{Stakeholder}:] Según el~\cite{PMI2019}, un \textit{stakeholder} es todo aquel individuo, grupo u organización que puede afectar, verse afectado por, o percibir que es afectado por una decisión, actividad o resultado de un proyecto, programa o portafolio.
	
	\item[Seguridad de la información:] Preservación de la confidencialidad, integridad y disponibilidad de la información, así como de otras propiedades asociadas como autenticidad, trazabilidad, responsabilidad y confiabilidad~\cite{ISO27000Glossary}.
	
	\item[SIEM (Security Information and Event Management):] Plataforma orientada a la recopilación, centralización, correlación y análisis de eventos y registros de seguridad provenientes de múltiples dispositivos, sistemas y servicios dentro de una infraestructura tecnológica~\cite{IBM_SIEM}.
\end{description}


\section*{T}
\begin{description}
	\item [Trazabilidad:] Capacidad de asociar de manera inequívoca las acciones, eventos o actividades realizadas sobre la información o un sistema a un usuario, proceso o entidad específica~\cite{ISO27000Glossary}.
\end{description}


\section*{V}
\begin{description}
	\item[Virtualización:] Tecnología que permite crear versiones virtuales de recursos computacionales, tales como servidores, sistemas operativos, redes o dispositivos de almacenamiento, sobre una misma infraestructura física~\cite{NISTVirtualization}.
	
	\item[Vulnerabilidad:] Debilidad presente en un activo, sistema o control de seguridad que puede ser explotada por una o más amenazas~\cite{ISO27000Glossary}.
\end{description}

\section*{Z}
\begin{description}
	\item[Zero-Trust Architecture (ZTA):] Modelo de seguridad basado en el principio de “nunca confiar, siempre verificar”, en el cual todo acceso a recursos, sistemas o servicios debe ser autenticado, autorizado y validado continuamente~\cite{NIST_ZTA}.
\end{description}
